\documentclass{article}
\usepackage{graphicx} % Required for inserting images
\usepackage[a4paper,margin=2.2cm]{geometry}
\usepackage{amsmath}
\usepackage{booktabs}
\usepackage{float}
\usepackage[font=small,skip=4pt]{caption}
\usepackage{hyperref}

\newcommand{\repo}{\url{https://github.com/richbang/ELEC872-Lab1}}

\title{ELEC 872 Week4 Workshop}
\author{Daeho Bang, Yujong Cho}
\date{October 6, 2026}

\begin{document}

\maketitle

\section{Baseline result}

LBP$^{u2}_{16,2}$, $30\times37$ windows, $\chi^2$ distance, nearest neighbour on ORL~\cite{ahonen2004face}.
Images 1--5 of each subject are the gallery and images 6--10 the probes (200 each); only the
probes are changed (uneven lighting: $I + 255\,s\,(x-0.5)$; rotation: about the image centre).
Our run reproduces the numbers given in the workshop description.

\begin{table}[H]
\centering
\begin{tabular}{lccc}
\toprule
 & normal & light $s=0.8$ & rot $20^\circ$ \\
\midrule
Workshop description & --    & 0.705 & 0.630 \\
Our baseline run     & 0.975 & 0.705 & 0.630 \\
\bottomrule
\end{tabular}
\caption{Rank-1 recognition rate of the baseline.}
\end{table}

\section{Our improvement and result}

The LBP descriptor and $\chi^2$ matching are unchanged. We add two steps:
\begin{enumerate}
  \item \textbf{Brightness-plane subtraction}~\cite{sung1998example}: fit a plane
        $ax+by+c$ to the image by least squares and subtract it, which removes the
        left-to-right lighting gradient.
  \item \textbf{Rotation search}: store each gallery image rotated by
        $-24^\circ, -18^\circ, \dots, 24^\circ$ and use the smallest $\chi^2$ over these copies.
\end{enumerate}

\begin{table}[H]
\centering
\begin{tabular}{lccccccccc}
\toprule
 & & & \multicolumn{4}{c}{light $s$} & \multicolumn{3}{c}{rotation} \\
\cmidrule(lr){4-7}\cmidrule(lr){8-10}
 & normal & bright & 0.2 & 0.4 & 0.6 & 0.8 & $5^\circ$ & $10^\circ$ & $20^\circ$ \\
\midrule
Baseline & 0.975 & 0.975 & 0.955 & 0.940 & 0.870 & 0.705 & 0.965 & 0.890 & 0.630 \\
Ours     & \textbf{0.995} & \textbf{0.980} & \textbf{0.985} & \textbf{0.990} & \textbf{0.985} & \textbf{0.985} & \textbf{0.995} & \textbf{0.980} & \textbf{0.940} \\
\bottomrule
\end{tabular}
\caption{Baseline vs.\ ours on the workshop conditions.}
\end{table}

\noindent Accuracy on normal photos is not lost: over 100 random 5+5 splits (the paper's
protocol) the baseline gives 0.983 and ours 0.983. Ours also improves conditions we never
tuned on, e.g.\ rotation by $30^\circ$ ($0.345 \to 0.900$).

\section{Code}

\noindent Repository: \repo\\
Method: \texttt{lbpface/robust.py} \quad Evaluation: \texttt{experiments/workshop.py}\\
Run: \texttt{python -m experiments.workshop}

\bibliographystyle{plain}
\bibliography{references}

\end{document}
